\section{Introduction}
A hardware resource has value only through the programs that use it. Larger tiles can reduce repeated weight transfers, yet increase register demand enough to reduce occupancy. More local storage can enable reuse, yet consume area that could support additional execution units or outstanding memory requests. Fusing two operators can remove intermediate transfers, yet change the instruction mix and the resources that limit progress. These choices jointly determine the efficiency of AI infrastructure.

The interactions are visible even within a small design family. In a Ventus study, increasing a matrix-multiplication tile from $32\times64$ to $64\times64$ reduced repeated weight traffic, but the new kernel required 173 vector registers. It performed poorly on the preceding register configuration and became useful after the hardware allocation changed. An attention mapping rejected on the default machine subsequently improved execution on another configuration (Section~\ref{sec:tiling-results}). Such reversals expose a limitation of committing to software or hardware choices before their interactions are understood.

Layered design makes a large problem manageable. A computation graph describes operators, a compiler constructs programs, an instruction interface specifies executable operations, and a hardware implementation supplies resources. Each layer hides details that other layers need not routinely consider. However, a hidden detail can become a consequential design variable. A coarse attention graph specifies matrix multiplication and normalization without exposing whether the score matrix is materialized, which tiles remain live, or how partial normalization state is updated. Once these choices become explicit, the available algorithms and the hardware needed to run them can change together.

This leads to two connected difficulties. \emph{Representation} determines which designs can be considered. An optimizer cannot select an online reduction, a different communication pattern, or a specialized data path that its representation does not admit. \emph{Search} determines which of the admitted designs can be found within the available effort. Expressing all relevant tiling, storage, scheduling, instruction, and hardware choices at once can produce a mixed-integer program (MIP) with an impractical number of variables and weak relaxations. Improving only one side leaves the other as a bottleneck.

We organize this process as Agentic Progressive Co-Design. Agents construct and revise the mathematical design problem, while optimization solvers search its explicit choices. The overall domain includes algorithmic organization, executable programs, the software--hardware boundary, and hardware structure. The active formulation exposes a selected part of that domain. Agents can introduce a new algorithmic family, expand a resource abstraction, derive a stronger encoding, or reopen a group of coupled decisions. They retain the assumptions behind reductions so that an unfavorable result on one machine does not eliminate a useful implementation elsewhere.

Three foundations suggest how to organize this process. CoSA shows how regular computation and architectural constraints support MIP-based scheduling~\cite{cosa}. Axon demonstrates semantics-driven generation of tensor-program alternatives and instruction implementations~\cite{axon}. LLM4MIP demonstrates a complementary role for language models in analyzing difficult optimization instances and developing primal and dual strategies\footnote{LLM4MIP released workflows: \url{https://llm4mip.github.io/}.}. Our framework connects these directions around a changing software--hardware design domain: the program family, hardware capabilities, and optimization representation can all be revised.

The software--hardware boundary is an explicit part of that domain. A transfer pipeline may be orchestrated by software, exposed through programmable asynchronous instructions, or managed by a specialized controller. These alternatives change instruction overhead, queue requirements, control area, and the range of workloads supported. Choosing among them requires the same joint treatment as choosing tile dimensions and memory capacity. The objective is a matched implementation: an algorithm and program organization, an execution interface, and a hardware structure that together satisfy the workload and deployment constraints.

The paper develops this design process in three stages. Section~\ref{sec:space} defines the relationship between the bounded implementation domain and the designs represented by an active MIP. This distinction makes representation expansion, hardware-dependent dominance, and branch retention precise. Sections~\ref{sec:construction} and~\ref{sec:search} turn the formulation into a search procedure: tensor dependencies connect algorithmic alternatives to resource constraints, boundary variables price specialization, and solver and execution feedback guide subsequent model changes. Section~\ref{sec:evaluation} then examines these choices through a programmable architecture benchmark and Ventus experiments, using performance gains, preference reversals, storage tradeoffs, and model revisions to explain the value of joint design.

The empirical studies exercise executable software alternatives, parameterized hardware models, and MIP-guided selection. They establish the interactions that a progressive formulation must capture. Autonomous construction of the full hierarchy and generation of new instruction interfaces extend beyond the evaluated system; the formulation identifies how to incorporate them. Functional acceptance remains a prerequisite for comparing candidate designs.
